\section{Sharp integer degree bounds}\label{sec:integer-envelopes}

For integers \(N\ge1\) and \(R\ge0\), write \(R=aN+w\) with
\(0\le w<N\), and define
\begin{equation}\label{eq:phi-min}
 \Phimin(N,R)=(N-w)\binomtwo a+w\binomtwo{a+1}.
\end{equation}
If \(L\le U\), \(NL\le R\le NU\), put
\(E=R-NL=t(U-L)+w\), where \(0\le w<U-L\), and define
\begin{equation}\label{eq:phi-max}
 \Phimax(N,L,U,R)
 =N\binomtwo L+LE+t\binomtwo{U-L}+\binomtwo w.
\end{equation}
When \(L=U\), the latter expression has its evident constant interpretation.

\begin{lemma}[Balanced and filled degree sequences]\label{lem:degree-extrema}
\coverage{fragment}
Among integer sequences \(z_1,\dots,z_N\) with sum \(R\), the minimum of
\(\sum_i\binomtwo{z_i}\) is \(\Phimin(N,R)\).  With the additional bounds
\(L\le z_i\le U\), its maximum is \(\Phimax(N,L,U,R)\).
\lean{balanced_twicePairCount_lower}
\end{lemma}

\begin{proof}
If \(z_i\ge z_j+2\), replacing \((z_i,z_j)\) by
\((z_i-1,z_j+1)\) lowers the sum by \(z_i-z_j-1\).  Repetition leaves only
the adjacent values \(a,a+1\), proving \eqref{eq:phi-min}.  For the maximum,
the reverse exchange moves one unit from a smaller nonminimal coordinate to a
larger nonmaximal one and raises the sum.  Hence all but at most one coordinate
lie at \(L\) or \(U\); counting the filled coordinates gives
\eqref{eq:phi-max}.
\end{proof}

\begin{theorem}[Exact integer incidence condition]\label{thm:paired-envelope}
\coverage{fragment}
Let \(\mathcal D\) be a symmetric \(2\)-\((v,K,\Lambda)\) design whose point
set is partitioned into \(A\sqcup B\).  Let \(\cS\) be a family of \(T\)
blocks, each meeting \(A\) in \(s\) points.  Suppose the selected-block
degrees satisfy integer bounds
\[
 L_x\le z_x\le U_x\qquad(x\in A\sqcup B).
\]
At the fixed totals
\[
 \sum_{x\in A}z_x=sT,
 \qquad
 \sum_{x\in B}z_x=(K-s)T,
\]
let \([P_A^-,P_A^+]\) and \([P_B^-,P_B^+]\) be the exact minimum and maximum
of the two sums \(\sum\binomtwo{z_x}\) over their respective integer boxes.
Then
\begin{equation}\label{eq:interval-overlap}
 [P_B^-,P_B^+]
 \cap
 \left[\Lambda\binomtwo T-P_A^+,
       \Lambda\binomtwo T-P_A^-\right]\ne\varnothing.
\end{equation}
For uniform boxes the endpoints are given by
Lemma~\ref{lem:degree-extrema}; heterogeneous endpoints are obtained by
integer water filling.
\lean{integerPairIntervals_overlap}
\end{theorem}

\begin{proof}
Count triples \((x,S_1,S_2)\) with \(S_1,S_2\in\cS\), \(S_1\ne S_2\), and
\(x\in S_1\cap S_2\).  Two blocks of a symmetric design meet in exactly
\(\Lambda\) points, so
\begin{equation}\label{eq:pair-partition}
 \sum_{x\in A}\binomtwo{z_x}
 +\sum_{x\in B}\binomtwo{z_x}
 =\Lambda\binomtwo T.
\end{equation}
Each sum lies in its exact integer interval.  Reflecting the interval for
\(A\) through the right side of \eqref{eq:pair-partition} gives
\eqref{eq:interval-overlap}.
\uses{lem:degree-extrema}
\end{proof}

The real relaxation is familiar, but the following formulation makes the
loss of information exact.

\begin{corollary}[Incidence mixing as the real relaxation]
\label{cor:mixing-relaxation}
\coverage{absent}
In the setting of Theorem~\ref{thm:paired-envelope},
\begin{equation}\label{eq:real-relaxation}
 T\left(\frac{s^2}{|A|}+\frac{(K-s)^2}{|B|}-\Lambda\right)
 \le K-\Lambda.
\end{equation}
This is algebraically identical to the incidence-graph mixing inequality for
the point set \(B\) and block family \(\cS\).  Equality requires constant
selected-block degree on both point classes.
\imports{LundSaraf:incidence-bound,MurphyPetridis:variance-identity}
\end{corollary}

\begin{proof}
Cauchy--Schwarz and \eqref{eq:pair-partition} give
\[
 \frac{s^2T^2}{|A|}-sT
 +\frac{(K-s)^2T^2}{|B|}-(K-s)T
 \le \Lambda T(T-1).
\]
Divide by \(T\) and use \(K=s+(K-s)\).  Conversely, expanding the standard
mixing inequality and substituting
\(K(K-1)=\Lambda(v-1)\) gives the same expression.  Equality in the two
Cauchy inequalities is equivalent to constant degrees.
\uses{thm:paired-envelope}
\end{proof}

For uniform lower bounds without an upper cap, the balancing remainder has the
closed form
\begin{equation}\label{eq:rounding-remainder}
 \Phimin(N,R)=\frac{R^2/N-R}{2}+\frac{w(N-w)}{2N},
 \qquad R\equiv w\pmod N,\qquad 0\le w<N.
\end{equation}
The last term is precisely what the real relaxation discards.  Along some
rational equality families, it has linear size and cannot be absorbed into
lower-order error.
